This project explores the development and feasibility of a gait recognition-based authentication system for physical access control, aiming to provide a less obtrusive and more user-independent alternative to traditional biometric and token-based systems.
By leveraging \ac{cv} and \ac{ml} techniques, the system identifies individuals based on their lateral walking patterns, enabling passive authentication without the need for direct user interaction.
The project proposes and implements a prototype that uses a client-server architecture, where the client captures gait data and the server performs identification and authorisation checks.
Through functionality, performance, and user testing, the prototype demonstrates that gait recognition can serve as a viable biometric authentication method, offering advantages in usability and transparency while maintaining an acceptable level of security.
The results suggest that with further refinement, gait-based authentication could enhance modern access control systems, particularly in environments where convenience, hygiene, or unobtrusiveness is prioritised.